\section{A deterministic recursion on width and mass}
\label{sec:algorithm}

We now give the procedure whose output meets the simultaneous
obligations of Theorem~\ref{thm:labelling}. Every step uses exact
integers or rationals and label vectors from children. The witnesses
from the preceding section are proof objects; no execution searches
for or tests them.

For an integer pivot $k$ and an even width $r$, a \emph{centered
call} with masses $(a',b')$ means
$\Label(k-r/2,r,a',b')$, on the box $[k-r/2,k+r/2]$.
This notation will only be used with integer half-widths.

In each pivot pass, one side's masses remain unchanged. Its translated
witnesses force labels at all gaps at most $p=q/32$; those labels will
keep these coordinates fixed while the pivot moves. The later
width-$D/2$ call has unchanged masses on both sides and makes all gaps
below $q$ into claims, supplying the labels for reweighting.

\begin{procedure}[$\Label(A,D,a,b)$]
\label{alg:label}
Set $u=A+D$ and $P=D/32$. Execute the first applicable base case;
if neither applies, proceed through Steps~3--6.
\begin{enumerate}[label=\textbf{\arabic*.},leftmargin=*]
\item\label{step:mass-base}\textbf{Mass base case.}
If $a_i b_i>1$ for all $i\in V$, return $+$ where $b_i>1$ and
return $-$ at the remaining vertices.

\item\label{step:width-base}\textbf{Width base case.}
If $D<4096$, compute the box solution $z$ by $\Basic$.
Return $+$ at $i$ if $z_i>A_i+D/2$, and $-$ otherwise.

\item\label{step:down}\textbf{Downward pass.}
Use $c,M,q,p$ from Equation~\eqref{eq:pivot-parameters} and
initialize $k=c$. Repeatedly make a centered width-$q$ call with
masses $(a,\frac43 b)$ and, from its labels, set simultaneously
\[
 k_i^{\mathrm{new}}=
 \begin{cases}
   \max\{c_i-M,k_i-p\},&\text{if the label at }i\text{ is }-,\\
   k_i,&\text{if the label at }i\text{ is }+.
 \end{cases}
\]
Replace $k$ by $k^{\mathrm{new}}$. Stop when the update leaves
$k$ unchanged, denote this last pivot by $k^{\downarrow}$, and
record
\[
                       I_- =\{i:k_i^{\downarrow}<c_i-P\}.
\]

\item\label{step:up}\textbf{Upward pass.}
Start from $k=k^{\downarrow}$. Repeatedly make a centered width-$q$
call with masses $(\frac43 a,b)$ and set simultaneously
\[
 k_i^{\mathrm{new}}=
 \begin{cases}
   \min\{c_i+M,k_i+p\},&\text{if the label at }i\text{ is }+,\\
   k_i,&\text{if the label at }i\text{ is }-.
 \end{cases}
\]
Replace $k$ by $k^{\mathrm{new}}$. Stop when it is unchanged,
write $k^f$ for the final pivot, and record
\[
                       I_+=\{i:k_i^f>c_i+P\}.
\]

\item\label{step:provisional}\textbf{Provisional labels and new masses.}
Make one centered width-$D/2$ call at $k^f$ with masses $(a,b)$.
For its returned vector $\lambda^0\in\{+,-\}^V$, define
\begin{equation}
 (\bar a_i,\bar b_i)=
 \begin{cases}
  (\frac45 a_i,\frac85 b_i),&\lambda_i^0=+,\\[1mm]
  (\frac85 a_i,\frac45 b_i),&\lambda_i^0=-.
 \end{cases}
 \label{eq:reweighting}
\end{equation}

\item\label{step:continuation}\textbf{Continuation and overrides.}
Call $\Label(A,D,\bar a,\bar b)$. In that child's output put
$-$ on $I_-$ and $+$ on $I_+\setminus I_-$, leaving every
other label unchanged. Return the resulting vector.
\end{enumerate}
\end{procedure}

Each pass includes a final child call whose update changes no
coordinate. Correctness will use that call's guarantee together with
the stationary update. The two passes run in the stated order;
within one update every coordinate uses the same returned label vector
and the same old pivot.

\subsection{Integral nested boxes and bounded passes}

The next facts depend only on the mechanics of the procedure, not on
the meaning of the returned labels.

\begin{lemma}
\label{lem:well-formed}
Every recursive child is a valid input of the form
Equation~\eqref{eq:recursive-input}, and its box is contained in
its parent's box. In a nonbase call the pivot satisfies
\[
 \begin{aligned}
 c-M\le k\le c&\qquad\text{during the downward pass},\\
 c-M\le k\le c+M&\qquad\text{during the upward pass}.
 \end{aligned}
\]
If child calls terminate, each pass makes at most $1024n+1$ calls.
\end{lemma}

\begin{proof}
In the recursive case, $D\ge4096$ is dyadic, so all parameters
in Equation~\eqref{eq:pivot-parameters}, both child half-widths,
and every pivot coordinate are integral. The child widths $D/128$
and $D/2$ are positive powers of two. The update caps imply the
stated pivot bounds. A centered child has width $r\le D/2$,
and its endpoints satisfy
\[
 c-M-r/2\ge A+D/8\ge A,
 \qquad c+M+r/2\le A+7D/8\le A+D.
\]
The continuation has the unchanged parent box. All mass multipliers
preserve positivity and rationality.

The pivot lies on the grid of spacing $p$ through $c$, since
$M/p=512$ is integral and both caps belong to that grid. In a
changing update at least one coordinate moves by at least $p$.
Each pass is coordinatewise monotone over a range of at most $2M$,
so a coordinate moves at most $2M/p=1024$ times. At most $1024n$
updates can change the pivot. One further child call certifies the
unchanged update at which the pass stops.
\end{proof}

\subsection{A rank that decreases for every child}

For positive masses define
\begin{equation}
 \mu(a,b)=\min_{i\in V}a_i b_i,
 \qquad \rho=\frac{32}{25},
 \qquad
 B(\mu)=
 \begin{cases}
  0,&\mu>1,\\
  1+\lfloor\log_{\rho}(1/\mu)\rfloor,&0<\mu\le1.
 \end{cases}
 \label{eq:budget}
\end{equation}
The integer $B$ measures how many guaranteed increases of the
minimum mass product can precede the mass base case. It is used for
analysis only; the executable procedure does not evaluate logarithms.

\begin{lemma}[Well-founded recursion]
\label{lem:termination}
The procedure terminates on every input in
Equation~\eqref{eq:recursive-input}. Writing $D=2^d$, every
child of a nonbase call has smaller rank
\[
                            R=d+B(\mu(a,b)).
\]
The provisional child in Step~5 is the only child that preserves $B$.
Every other child decreases $B$ by at least one.
\end{lemma}

\begin{proof}
If $\mu>1$, the mass base case applies. Otherwise multiplying
$\mu$ by a factor at least $\rho$ decreases $B$ by at least
one: either the new value exceeds one, or its logarithm in
Equation~\eqref{eq:budget} falls by at least one before rounding
down.

In either pivot pass, each product $a_i b_i$ is multiplied by
$4/3\ge\rho$. In the continuation, each product is multiplied by
$(4/5)(8/5)=\rho$, independent of the provisional label. Thus
these children decrease $B$ and do not increase $d$. The
provisional child keeps the masses and halves the width, decreasing
$d$ by one. Every child has strictly smaller nonnegative rank $R$.

Induction on $R$ now proves termination. Base cases require only
finite arithmetic and, in the width case, the finite iteration of
Lemma~\ref{lem:accepted-box}. At a nonbase node, every child
terminates by induction. Lemma~\ref{lem:well-formed} then bounds
the two passes; the remaining two children and local operations are
finite as well. This reasoning makes no correctness assumption about
a label.
\end{proof}

The distinction between width reduction and budget reduction is shown
in Figure~\ref{fig:recursion}. It permits induction for correctness
and a sharper count than a general branching bound: only one child
can preserve the mass budget, irrespective of the width.

\begin{figure}[htbp]
 \centering
 \begin{tikzpicture}[
 >=Stealth,
 box/.style={draw=linkblue,rounded corners=2pt,align=center,
             font=\small,inner sep=6pt,text width=3.6cm},
 edge/.style={->,draw=linkblue,thick}]
 \node[box,text width=4.5cm] (root) at (0,0)
   {Parent input\\width $D$, mass budget $B$};
 \node[box] (loops) at (-4.4,-2.1)
   {Pivot-pass children\\$\le 2(1024n+1)$ calls\\
    width $D/128$\\budget $\le B-1$};
 \node[box] (middle) at (0,-2.1)
   {Provisional child\\One call\\width $D/2$\\budget $B$};
 \node[box] (last) at (4.4,-2.1)
   {Continuation\\One call\\width $D$\\budget $\le B-1$};
 \draw[edge] (root.south west) -- (loops.north);
 \draw[edge] (root.south) -- (middle.north);
 \draw[edge] (root.south east) -- (last.north);
\end{tikzpicture}

 \caption{Children of a nonbase call. The two passes together produce
 at most $2(1024n+1)$ narrow children, each decreasing $B$.
 The provisional child preserves $B$ and halves the width; the
 continuation preserves the box and decreases $B$. Edges show the
 recursion tree, not the execution order.}
 \label{fig:recursion}
\end{figure}
